\section{A smaller denominator also shortens the compact equations}
\label{sec:xy-carrier}

The direct-calculus quartic and the compact positive-integer equations serve
different purposes. This section retains the existing native/K215 calibration
and its fixed coefficient tuple, but makes a further proved reduction. It
gives ordinary polynomials in ten strictly positive integers with
operation/degree pairs
\[
 (77,626),\qquad(78,464),\qquad(79,460).
\]
The first two improve both metrics over the immediate $(78,630)$ and
$(79,468)$ parents. The third improves the second parent's degree while
preserving its operation count. These are representation tradeoffs; they
are not claims of global minimality.

\subsection{Complete polynomial}

The only unknowns are
\[
 J,F,\alpha,i,T,w,s,n,\eta,\zeta\in\mathbb Z_{>0}.
\]
The four fixed coefficients $\beta,K_0,M_C,M_F$ are the complete directly
indexed BFS943 tuple retained in the archive. Their exact little-endian
bytes and the full finite reconstruction recipe are supplied. The largest
coefficient port has $592{,}884$ bits. These ports are fixed numerals, not
additional unknowns. They satisfy
\[
 \beta\ge31\text{ odd},\quad K_0>0,\quad
 0<M_C<\beta<M_F<2\beta.
\]
All following symbols are polynomial abbreviations:
\begin{gather*}
 q=\beta J+1,\quad h=q-1,\quad C=q-F-\alpha,
 \quad X=wq,\quad Y=q^3s,\\
 k=\eta+\zeta,\quad c=kY+\eta,\quad
 A=Y(X+1)+1,\quad\Delta=A^2-1,\quad H=4A-5,\\
 V_0=XY^2,\qquad E=1+V_0(V_0+1)k^2,\\
 R=[h(q-F)+\alpha](q^2-1)+(M_C+qM_F)J,\\
 e=k-R-1,\quad B=(K_0+w)C-F,\quad U=ic^2,\\
 a=\Delta U^2E,\quad F_0=1+a,\quad G=EF_0,\\
 F_*=1+4aF_0,\quad Q=\Delta(F_*-1),\\
 b=c(T+\alpha)-R,\quad V=F_*b-c,\quad
 K_+=X(3c-2X)+2.
\end{gather*}
Choose $(D,L)=(XY,V^2)$, $(XY,bV)$, or $(q^4,bV)$, respectively,
and put
\begin{align*}
 M&=hDH,\\
 Z&=hD\bigl[H(n-LB)+K_+\bigr]-H(DB+he),\\
 R_G&=M^2G.
\end{align*}
The single equation is
\begin{equation}
 \boxed{P_{D,L}=Q(Z^2-R_GV^2)+(R_G-Z^2)=0.}
 \label{eq:xy-carrier}
\end{equation}
There are no extra quantified variables, square-root functions, divisions,
or variable exponents in this equation. Its three literal arithmetic
circuits are included as JSON.

\subsection{The admissible-denominator argument}

The new lemma applies to any specified polynomial $D=D(q,X,Y)$ satisfying
\begin{equation}
 D>2,\qquad q^4\mid D,\qquad D\mid4XY^2
 \label{eq:xy-admissible}
\end{equation}
on the positive inputs. Both chosen denominators satisfy these conditions:
$XY=q^4ws$, and $4XY^2/q^4=4wq^3s^2$.

For proof notation only define
\[
 z=Z/M,\quad \kappa=K_+/H,\quad\rho=e/D,\quad
 t=B/h,\quad W=1+hL.
\]
Identically,
\[
 P_{D,L}=(Q-1)Z^2-M^2G(QV^2-1),\qquad
 z=n+\kappa-Wt-\rho.
\]
The unchanged product-square argument uses
$F_*=(1+2a)^2$ and $Q=4\Delta^2U^2G$ to extract
$G=g^2$, $E=\tau^2$, and $F_0=f_0^2$ with positive integral roots.
Adjacent-Pell descent then proves
\[
 QV^2-(Q-1)y^2=1,\qquad z=\varepsilon gy\in\mathbb Z.
\]
The unchanged doubling and folding give a strong pair with $c^2\mid I$
and a main index $c=\psi_A(p)$. The first norm gives
\[
 k=2\psi_{P_0}(n_0),\qquad P_0=2XY^2+1.
\]
The retained sign-transport estimates require only $D>2$; they prove
$B\ge0$, $C>0$, and then identify $p=R$ with $3q<R<q^4-2$.
No integrality of $\rho$ has been assumed in these steps.

Now the integral carrier, after clearing $hH$, gives $D\mid hHe$.
Since $q^4\mid D$, $h=q-1$, and $H\equiv-1\pmod q$, it follows that
$q^4\mid e$. The base-one Pell recurrence gives
\[
 k\equiv2n_0\pmod{4XY^2},
 \qquad q^4\mid 2n_0-R-1.
\]
Also $P_0>A$. If $n_0\ge R$, monotonicity would imply
$k/2\ge\psi_A(R)=c$, contradicting $c>k$. Hence $1\le n_0<R$, so
\[
 |2n_0-R-1|<R<q^4.
\]
Thus $2n_0=R+1$. Returning to the full Pell congruence proves
$4XY^2\mid e$, and \eqref{eq:xy-admissible} now proves $D\mid e$.
The full denominator is recovered only after the index. No coprimality
claim about all factors of $D$ is required.

With $\rho$ integral, the inherited projection modulo $H$ forces
$X=2^R$, and then $\kappa\in\mathbb Z$ and $h\mid B$. The unchanged
parity and half-binomial reconstruction recover the native source and
therefore a finite ZFC refutation.

Conversely, choose the inherited canonical witnesses from a refutation.
They have positive $e$ divisible by $4XY^2$, so $\rho=e/D$ is a positive
integer for either new denominator. Keep the fresh strong and auxiliary
Pell witnesses and set
\[
 n=gy+(1+hL)t+\rho-\kappa.
\]
All terms are integral, and the retained estimate $g>\kappa$ proves $n>0$.
The carrier equals $gy$, giving a zero of \eqref{eq:xy-carrier}. The
dependence $D=D(q,X,Y)$ ensures that choosing $n$ last is noncircular.
This proves both existential directions and hence, by negation,
\[
 \mathrm{ZFC}\vdash
 \bigl(\neg\exists\mathbf x\in\mathbb Z_{>0}^{10}\,
 P_{D,L}(\mathbf x)=0\bigr)
 \longleftrightarrow\operatorname{Con}(\mathrm{ZFC}).
\]
The full ordinary proof and explicitly identified inherited lemmas are
retained in the compact packet.

\subsection{Exact cost and degree}

For $D=XY$ the three operations
\[
 D=XY,\qquad A-1=D+Y,\qquad V_0=DY
\]
replace four parent operations. The $q^4$ carrier adds one square of the
already live value $q^2$. The resulting exact counts are
\begin{center}
\begin{tabular}{lrrrrr}
\hline
$(D,L)$&Multiplications&Additions&Subtractions&Total&Degree\\
\hline
$(XY,V^2)$&45&18&14&77&626\\
$(XY,bV)$&46&18&14&78&464\\
$(q^4,bV)$&47&18&14&79&460\\
\hline
\end{tabular}
\end{center}
Writing $q_0=\beta J$, $k_0=\eta+\zeta$, and $t_0=T+\alpha$, the full
leading homogeneous forms are
\begin{align*}
 [P_{XY,V^2}]_{626}
 &=2^{14}q_0^{378}w^{48}s^{110}i^{20}k_0^{64}t_0^4
       (q_0-F-\alpha)^2,\\
 [P_{XY,bV}]_{464}
 &=-2^{10}q_0^{284}w^{36}s^{82}i^{14}k_0^{46}t_0^2,\\
 [P_{q^4,bV}]_{460}
 &=-2^{10}q_0^{284}w^{34}s^{80}i^{14}k_0^{46}t_0^2.
\end{align*}
These are nonzero for the actual fixed positive $\beta$. The independent
checker verifies every coefficient of these forms and 28 exact polynomial
identities covering every circuit gate. The actual fixed coefficients are
reconstructed from their finite recipe, and altered gates are rejected.
Such finite checks establish the delivered algebra; the unbounded
consistency equivalence is the ordinary proof above and its stated
inherited calibration.
